\documentclass[../main]{subfiles}
\begin{document}
\ifSubfilesClassLoaded{\setcounter{page}{50}}{}
\section{Комунізм і сім’я}
\label{sec:09_kommunizm_semya}

У праці «Походження сім’ї, приватної власності та держави» Фрідріх Енгельс писав: «Згідно з матеріалістичним розумінням, визначальним моментом в історії є, зрештою, виробництво та відтворення безпосереднього життя. Але саме воно, знову ж таки, буває двох видів. З одного боку – виробництво засобів до існування: продуктів харчування, одягу, житла та необхідних для цього знарядь; з іншого – виробництво самої людини, продовження роду. Суспільні порядки, за яких живуть люди певної історичної епохи та певної країни, зумовлені обома видами виробництва: ступенем розвитку, з одного боку, – праці, з іншого – сім’ї».

У попередніх розділах ми говорили про те, що за комунізму виробництво життєво необхідних предметів автоматизовано та підпорядковано виробництву людського життя в усій його повноті. Суспільство свідомо творить свій власний розвиток, свої власні форми колективності, розширюючи зміст вільної діяльності. Сім’я, як господарська одиниця, заснована на приватному характері присвоєння продуктів праці, втрачає своє значення та зникає в перехідний період. «Приватне домашнє господарство перетвориться на суспільну галузь праці», – писав Енгельс. Вище ми бачили, наскільки це заощаджує суспільний час, звільняючи людей від тягарів побуту. Однак об’єднання стосується не лише споживання.

Найбільша революція у сфері сімейних відносин полягає в суспільному вихованні дітей. Ця революція давно назріла, і рух у цьому напрямку вже існує за капіталізму. Домашнє господарство почало «суспільно організовуватися», побут давно значною мірою став галуззю суспільного господарства. Ми готуємо їжу переважно з напівфабрикатів, купуємо одяг у магазині, перемо в пральні, хоча зовсім нещодавно це була жіноча справа, яка відбувалася вдома. Сучасному чоловікові також не потрібно носити воду, рубати дрова та будувати дім, в якому живе сім’я. Все це вже давно є суспільною високотехнологічною справою.

А от виховання дітей, чомусь, так і не перетворилося на сферу виробництва, хоча громадські заклади для цього з’явилися вже давно, і активно розвиваються не менше двох століть. І завжди ті, хто заможніший, віддавали дітей саме туди, наймали приватних викладачів, а самі лише насолоджувалися сімейним затишком. Так само і в Радянському Союзі – одним із перших досягнень, про яке з такою гордістю говорили соціалістичні трудящі, була організація ясел і дитячих садків при фабриках. Тобто потреба віддати виховання дітей тим, хто це вміє, вивільнити час для того, що вважаєш справді потрібним і важливим (для когось, до речі, це і є виховання дітей, чому б і ні) – у цьому єдині і бідні, і багаті, з деякими лише відмінностями міщансько-поетичного характеру. Капіталізм вже давно руйнує традиційну сім’ю, і цей процес не має жодних моральних підґрунть, а лише економічні. Зараз пропаганда сімейних цінностей зумовлена, з одного боку, досягненнями соціалістичних революцій, які заборонили експлуатацію дитячої праці, а з іншого – потребою капіталу у ринку дитячих товарів та освіти. Але справжньої історичної, суспільної необхідності у збереженні сім’ї в нинішній формі господарської одиниці не існує.

Отже, зі переходом суспільства до комуністичних відносин, сфера виховання дітей стане сферою суспільного виробництва. Виховання перетворюється на справжню соціальну творчість – вільну діяльність, що відповідає критеріям істини, добра та краси. Це значно підвищує відповідальність вихователів як перед дітьми, так і перед суспільством в цілому. Як колись писав А. С. Макаренко: «У нас не повинно бути жодних дитячих катастроф, жодних невдач, жодного відсотка браку, навіть вираженого сотими частками!» У багатьох працях він порівнював виховання з іншими галузями і показував, наскільки тут невдача небезпечніша, ніж деінде. Говорячи про сім’ю, Антон Семенович зазначав, що завдання батьків – виховати гідних членів суспільства, за «якість наданого матеріалу» вони несуть відповідальність перед суспільством, отримуючи натомість теплоту сімейних відносин. Але далеко не всі батьки вміють виховувати дітей. Крім того, далеко не всі діти мають батьків або повноцінну сім’ю вже зараз. У цьому полягає необхідність у свідомому, якісному розширенні можливостей виховання, у створенні інших форм «сім’ї».

Комуністична «сім’я» – це колективна форма, в якій кожен бере участь у створенні суспільного життя в суспільному масштабі. З цього стає зрозуміло, що розміри «сім’ї» можуть варіюватися, і, безумовно, вона не може обмежуватися сучасними рамками. Їх розмір залежатиме від багатьох конкретно-історичних умов суспільного виробництва, тому можливі різні варіанти як у межах однієї епохи розвитку комуністичного суспільства, так і в його історії. Оскільки така «сім’я» як первинний колектив, що формує особистість, базуватиметься на суспільній природі діяльності та безпосередньо визначатиметься нею, це означає, що принцип створення таких колективів також буде безпосередньо суспільним – якісно іншим, ніж зараз. Йдеться про знищення сучасного шлюбу як форми, тим більше, що статеві стосунки та статева любов і шлюбно-побутові відносини – це далеко не одне й те саме. Тому одним із найважливіших і незмінних практичних відкриттів А. С. Макаренка було те, що нова «сім’я» повинна безпосередньо включати кожну людину в суспільне виробництво, як активного члена суспільного трудового колективу, а не як споживача суспільних благ чи як «об’єкт» виховання.

Відстежуючи швидкі, а подеколи й непомітні неозброєним оком зміни в цій сфері вже в перші роки Великої Жовтневої соціалістичної революції, Олександра Коллонтай у своїй праці «Дорогу крилатому Еросу!» спробувала побачити та передбачити той тип особистих, точніше, відносин сексуальної любові, який утвердиться після повної перемоги суспільної власності над приватною, з розвитком нових суспільних інститутів і з впровадженням елементів інфраструктури, здатної полегшити тягар материнства, перетворюючи його на галузь суспільного виробництва.

Нове побутове життя, де воно було налагоджене (не лише в СРСР), справді звільнило жінок у незбагненних масштабах, про які інші країни могли лише мріяти, дало час не лише на саму любов та залицяння, а й на переосмислення любові. Такі нові, вже частково комуністичні відносини описані і в художній літературі часів індустріалізації, як у прозі самої Олександри Коллонтай («Василиса Малигіна», «Велика любов»), так і, наприклад, у Іллі Еренбурга в романі «День другий».

Запровадження комуни без реєстрації шлюбу (у 1920-1930-х роках не зареєстрований шлюб став нормою, а реєстрація – навпаки, ознакою недовіри між партнерами) там, де раніше панувала патріархальна сім’я, безперечно, було значним випередженням подій. Але це стало революцією у самосприйнятті тих, хто кохає, закоханих, членів сімей, розлучених. Не шлюб як життя для двох, заради двох, а сімейне життя, яке безпосередньо визначається цілями бурхливо зростаючого і розвиваючогося суспільства, найближчого (трудового) колективу – ось реальний зміст кохання, що виражає тенденцію до комунізму.

Сім’я як сфера любові та пов’язаної з нею відповідальності в капіталістичному суспільстві – це «зона підвищеного ризику». Подружжя приречене в межах власного житла переживати все те, що зумовлено відносинами приватної власності, питаннями достатку, заробітку, витрат та іншою буденністю життя. Але комунізм – це саме усунення таких проблем.

Динаміка зникнення як старої сім’ї, так і держави як такої є складною. Однак, часткові заходи соціалізму тут виявилися безсилими: з комун-гуртожитків студенти, молоді фахівці все ж опинялися у квартирах, у сім’ях, у старих рамках, але з новим соціальним забезпеченням. Мотиви для подолання старих форм «спілок» розгубилися в суспільстві, де грошово-товарні відносини знову відвойовували собі домінуюче становище. За поривом до комуни послідував закономірний і теж нібито тимчасовий поворот до патріархальності. Знову турбота лише батьків, а не всього колективу чи суспільства, виходила на перший план у вихованні. Адже намічалися не лише «сини полку» (війна вносила корективи), але й діти комун.

В цих «дітях комун» і втілювався один із найважливіших моментів для становлення комунізму. Під час переходу до комунізму діти повинні перестати ділитися на «своїх» і «чужих». Не має існувати «чужих» дітей (тобто дітей, про яких має піклуватися хтось інший, але не я) у величезному колективі, де знято проблему «підрахунку» особистого трудового внеску.

Потрапляючи в колектив, де багато чоловіків і жінок, діти, розрізняючи батька і матір, не будуть бачити «кінець світу» там, де закінчуються мама і тато і починаються «сусіди». Завдяки природній соціалізації в комуні, завдяки щоденному відчуттю турботи колективу, діти перестануть «грати в пісочниці приватної власності». Відмова від старого шлюбу дозволить уникнути низки пережитків, що відтворюються в психіці, у свідомості батьків і матерів, а не лише дітей. Зміна партнерства у великому та дружньому колективі перестає бути чимось травмуючим, якщо між членами пари вже немає любові. Взагалі, за відсутності суспільних економічних бар’єрів, усе суто особисте (в тому числі й інтимне) може стати справді особистим, але не таким, що протистоїть суспільству, звільнившись від господарських, правових і моральних пут.

Нова «сім’я» має сама стати безпосередньо колективною формою комуністичної діяльності, спрямованою на створення нової людини, безпосередньо зв’язуючи всіх її членів, включно з дітьми, із суспільством як дієвими суб’єктами.
\end{document}
